\documentclass[ 12pt,a4paper,twocolumn,fleqn]{article}
\usepackage{graphicx}
\usepackage[a4paper,top=20mm, bottom=30mm, left=10mm, right=25mm]{geometry}
\usepackage{fancyhdr}
\usepackage{hyperref}
\usepackage{listings}
\usepackage[]{algorithm2e}
\usepackage{color}
\usepackage{eso-pic}
\usepackage{tikz}
\pagestyle{fancy}

% ---- Automatic double-box border on EVERY page ----
\AddToShipoutPictureBG{%
  \begin{tikzpicture}[remember picture, overlay]
    \draw[line width=0.75pt]
      ([xshift=0.5cm, yshift=0.5cm]current page.south west) rectangle
      ([xshift=-0.5cm, yshift=-0.5cm]current page.north east);
    \draw[line width=0.75pt]
      ([xshift=0.7cm, yshift=0.7cm]current page.south west) rectangle
      ([xshift=-0.7cm, yshift=-0.7cm]current page.north east);
  \end{tikzpicture}
}
% ----------------------------------------------------

\usepackage{lineno}
\usepackage{xtab,booktabs}
\usepackage{setspace}
\usepackage{amsmath}
\usepackage{xcolor}
\usepackage{subcaption}
\usepackage[compact]{titlesec}
\titlespacing{\section}{0pt}{*0}{*0}
\titlespacing{\subsection}{0pt}{*0}{*0}
\titlespacing{\subsubsection}{0pt}{*0}{*0}
\setlength\columnsep{25pt}
\makeatletter
\g@addto@macro{\normalsize}{%
\setlength{\abovedisplayskip}{0pt}
\setlength{\abovedisplayshortskip}{0pt}
\setlength{\belowdisplayskip}{0pt}
\setlength{\belowdisplayshortskip}{0pt}}
\makeatother
\mathindent=0.0pt
\usepackage{float}
\renewcommand{\baselinestretch}{1.5}

\begin{document}
\lhead{CampusBite: University Canteen Management System}
\chead{}

\onecolumn
\begin{center}
\textcolor{brown}{\LARGE{INDEPENDENT UNIVERSITY, BANGLADESH}} \\
\includegraphics[scale=0.5]{IUB.png} \\

\large{Computer Science and Engineering} \\
\large{School of Engineering, Technology and Sciences} \\

\text{A Project Report On} \\
\smallskip
\textcolor{cyan}{\LARGE{CAMPUSBITE}} \\
\smallskip
\large{A Web-Based University Canteen Management System} \\



\textit{Submitted By} \\
\textbf{Sheikh Mayabee Hasan, \space 2430901} \\
\textbf{Aditya Kar, \space 2431020} \\

\textit{Under the guidance of} \\
\textbf{Md Masum Billah} \\
\text{Assistant Professor, CSE, IUB}\\
\large{\textbf{Summer: 2026}} \\
\textcolor{blue}{\large{Department of Computer Science and Engineering }} \\
\textcolor{blue}{\large{INDEPENDENT UNIVERSITY, BANGLADESH}} \\
\textcolor{blue}{\large{DHAKA}} \\
\end{center}

\newpage

\begin{center}
\LARGE{{\underline{Acknowledgement}}} \\
\end{center}
\normalsize

We would like to express our sincere gratitude to our supervisor, Md Masum Billah, Assistant Professor, Department of Computer Science and Engineering, Independent University, Bangladesh, for his continuous guidance, valuable feedback, and encouragement throughout the development of the CampusBite project. His support was instrumental in helping us shape the objectives, design, and implementation of this system.

We are also thankful to the Department of Computer Science and Engineering, IUB, for providing us with the opportunity, resources, and academic environment necessary to carry out this project. Finally, we would like to thank our families and friends for their constant support and motivation during the course of this work.

\begin{flushright}
\begin{centering}
    \textbf{Sheikh Mayabee Hasan \space 2430901} \\
\textbf{Aditya Kar \space 2431020} \\

\end{centering}

\end{flushright}


\newpage

\setstretch{1.5}
\twocolumn[
\begin{@twocolumnfalse}
\title {CampusBite: A Web-Based University Canteen Management System}
\date{}
\author{Sheikh Mayabee Hasan, Aditya Kar}
\maketitle



{\bfseries\LARGE Abstract} \\
    
University canteens are often difficult to navigate because students have no easy way to check what food is available, how much it costs, or whether a particular canteen currently has the items they want. At the same time, canteen owners struggle to keep menu information up to date and rarely receive structured feedback from students. CampusBite is a web-based university canteen management system designed to solve this problem by providing a centralized platform for both students and canteen owners. Students can register, log in, browse the canteens available on campus, view menus with prices and availability, and submit reviews and complaints. Canteen owners can log in to a dedicated dashboard to add, edit, and delete food items, update prices and availability, and view the feedback submitted by students. The system is built using HTML, CSS, and JavaScript for the front end, PHP for server-side processing, and MySQL for data storage, and is developed and tested using XAMPP and phpMyAdmin. This report presents the problem definition, related work, methodology, requirements, implementation, and results of the CampusBite system, and concludes with directions for future enhancement such as online ordering, payment integration, and a mobile application.


\end{@twocolumnfalse}]
\modulolinenumbers[5]
\onecolumn
\newpage
\tableofcontents
\newpage

\section{Introduction}

Canteens are an essential part of everyday university life, yet the process of finding food on campus is still largely informal. Students typically have to physically visit a canteen to find out what is being served, what it costs, and whether it is still available. If an item has run out or a canteen is closed, the student only discovers this after making the trip. There is also no structured way for students to give feedback: complaints about food quality, hygiene, or service are usually passed on verbally, if at all, and rarely reach the canteen owner in an organized way.

From the canteen owner's side, managing food information is done manually. Prices and availability are not communicated to students in real time, and there is no dashboard to track what students think about the food being served. This creates a communication gap between students and canteen owners that a simple web-based system can help close.

CampusBite is proposed as a solution to this gap. It is a web-based university canteen management system that gives students a single place to view all campus canteens, browse their menus, check prices and availability, and submit reviews and complaints. At the same time, it gives canteen owners a dashboard through which they can manage their own food listings and view the feedback submitted about their canteen. The system is supported by a MySQL database that keeps student, owner, canteen, food, review, and complaint information organized and consistent.

\subsection{Scope}

The scope of our CampusBite project is to develop a web-based university canteen management system that makes it easier for students to access canteen and food information. The system provides a centralized platform where students can view available campus canteens, explore food menus, check prices and food availability, and provide feedback through reviews and complaints.

The system also covers the needs of canteen owners. Owners can log in to their dashboard and manage their food information. They can add new food items, edit existing food details, delete food items, update prices, and change the availability status of food. They can also view reviews and complaints submitted by students.

\subsubsection{Main Areas Covered by the Project}

\textbf{1. Student Management}
\begin{itemize}
    \item Student registration and login
    \item Secure access to the student dashboard
    \item Viewing campus canteens
    \item Browsing food items
    \item Checking food prices and availability
    \item Submitting reviews
    \item Submitting complaints
\end{itemize}

\textbf{2. Canteen Owner Management}
\begin{itemize}
    \item Owner login
    \item Owner dashboard
    \item Adding food items
    \item Editing food information
    \item Deleting food items
    \item Updating food prices
    \item Updating food availability
    \item Viewing reviews and complaints
\end{itemize}

\textbf{3. Canteen and Food Management} \\
The system stores information about different campus canteens and the food they offer. Each food item is associated with a particular canteen. This allows students to find food according to the selected canteen.

\textbf{4. Review and Complaint Management} \\
CampusBite provides students with a structured way to give feedback. Students can submit reviews about canteens and raise complaints when they face problems. This information can be viewed by the relevant canteen owner.

\textbf{5. Database Management} \\
The project uses a MySQL database to store and organize the system's information. The main entities include:
\begin{itemize}
    \item Student
    \item Owner
    \item Canteen
    \item Food
    \item Review
    \item Complaint
\end{itemize}
These entities are connected through relationships to maintain organized and consistent data.

\textbf{6. Technology Scope} \\
The project is developed as a web application using:
\begin{itemize}
    \item HTML -- webpage structure
    \item CSS -- design and styling
    \item JavaScript -- client-side interaction
    \item PHP -- backend/server-side processing
    \item MySQL -- database management
    \item XAMPP/Apache -- local development server
    \item phpMyAdmin -- database management
\end{itemize}

\subsubsection{What is Outside the Current Scope}

The current version of CampusBite does not include:
\begin{itemize}
    \item Online food ordering
    \item Online payment
    \item Food delivery
    \item Food reservation
    \item Mobile application
    \item AI-based food recommendations
    \item Advanced inventory management
    \item Admin panel
\end{itemize}
These can be considered as future improvements.

\subsubsection{Future Scope}

In the future, CampusBite can be expanded by adding online food ordering and digital payment facilities. Students could reserve food, receive notifications about food availability, and use a mobile application. The system could also provide canteen ratings, real-time availability updates, and analytics to help owners understand popular food items and student preferences.

\subsubsection{In Simple Words}

The scope of CampusBite is basically:
\begin{itemize}
    \item \textbf{Students} $\rightarrow$ Find food + Check availability + Give feedback
    \item \textbf{Owners} $\rightarrow$ Manage food + Update availability + Handle feedback
    \item \textbf{Database} $\rightarrow$ Store all student, owner, canteen, food, review, and complaint information
\end{itemize}

\section{Problem Definition}

University students face several recurring problems when trying to access canteen services on campus:

\begin{itemize}
    \item Students have no way to check food prices or availability before physically visiting a canteen, which often leads to wasted time when an item is out of stock or the canteen is closed.
    \item There is no centralized platform that lists all the canteens on campus along with their respective menus, so students must rely on word of mouth or repeated visits.
    \item Students have no structured channel to submit reviews or complaints about food quality, hygiene, or service, which means genuine feedback rarely reaches the people who can act on it.
    \item Canteen owners have no digital tool to manage their menu, update prices, or mark items as unavailable, so this information is rarely kept current.
    \item Owners have no visibility into what students think about their canteen, since feedback is informal and undocumented.
\end{itemize}

CampusBite addresses these problems by providing a single web-based platform that connects students and canteen owners, gives students accurate and up-to-date food information, and gives owners the tools to manage their offerings and respond to feedback.

\section{Literature Review}

Canteen and cafeteria management is a recurring problem addressed in software engineering coursework and small-scale campus projects, since it combines multi-role access control, inventory-style data management, and user feedback in a single system. Several existing food-ordering and canteen management systems \cite{b1} demonstrate that a simple relational database design, built around entities such as user, vendor, item, and order, is sufficient to support the core operations of browsing, listing, and updating food information for a campus-scale deployment.

Studies on digital feedback systems \cite{b2} show that giving users a structured way to submit reviews and complaints, rather than relying on informal or verbal feedback, significantly increases the amount and quality of feedback that reaches service providers. This supports the design decision in CampusBite to treat reviews and complaints as first-class entities in the database rather than free-form, unstored text.

Work on PHP and MySQL-based web application architecture \cite{b3} highlights the common three-tier structure -- a browser-based front end, a PHP application layer, and a MySQL data layer -- as a practical and well-supported choice for small to medium web systems developed and deployed using tools such as XAMPP. This architecture was adopted for CampusBite because it is lightweight, widely documented, and appropriate for a system with a limited number of concurrent users, such as the students and canteen owners of a single university.

Taken together, the reviewed literature confirms that a role-based, database-driven web application is a suitable and proven approach for a university canteen management system, and it informed the methodology and system design decisions described in the following sections.

\section{Methodology}

The development of CampusBite followed a structured, incremental approach consisting of requirement gathering, system design, implementation, and testing.

\subsection{Requirement Gathering}

The requirements for CampusBite were gathered by identifying the day-to-day difficulties that students and canteen owners face on campus, such as the lack of visibility into canteen menus and the absence of a feedback channel. These observations were translated into functional and non-functional requirements, which are detailed in Section 5, and were used to define the two primary user roles supported by the system: \textit{Student} and \textit{Canteen Owner}.

\subsection{System Design}

Based on the gathered requirements, the system was designed around six core entities: Student, Owner, Canteen, Food, Review, and Complaint. A relational schema was designed in MySQL so that:
\begin{itemize}
    \item Each \textit{Canteen} is owned by one \textit{Owner}.
    \item Each \textit{Food} item belongs to exactly one \textit{Canteen}.
    \item Each \textit{Review} and \textit{Complaint} is submitted by a \textit{Student} and is associated with a specific \textit{Canteen}.
\end{itemize}
This design allows a student to browse food by canteen, and allows an owner to see only the food items, reviews, and complaints relevant to their own canteen.

\subsection{Implementation}

The system was implemented as a role-based web application with separate interfaces for students and canteen owners.

\subsubsection{Frontend Development}
The user interface was built using HTML for page structure, CSS for styling and layout, and JavaScript for client-side interactivity such as form validation and dynamic content updates. Separate views were created for the student dashboard (canteen listing, menu browsing, review and complaint submission) and the owner dashboard (food item management, price and availability updates, review and complaint viewing).

\subsubsection{Backend Development}
PHP was used to handle server-side logic, including user authentication for both students and owners, session management, and all create, read, update, and delete (CRUD) operations on canteens, food items, reviews, and complaints. PHP scripts communicate with the MySQL database to retrieve and persist data based on requests from the front end.

\subsubsection{Database Implementation}
The MySQL database was implemented and managed using phpMyAdmin, with tables corresponding to the Student, Owner, Canteen, Food, Review, and Complaint entities described in the system design. Foreign key relationships were used to enforce the associations between canteens and their owners, and between food items, reviews, and complaints and their respective canteens.

\subsection{Development Environment}
The system was developed and tested locally using XAMPP, which provides the Apache web server, PHP interpreter, and MySQL database required to run the application, along with phpMyAdmin for direct inspection and management of the database during development.

\section{Requirements}

\subsection{Functional Requirements}
\begin{itemize}
\item The system shall allow students to register and log in securely.
\item The system shall allow students to view a list of all available campus canteens.
\item The system shall allow students to browse the food items offered by a selected canteen, including price and availability status.
\item The system shall allow students to submit a review for a canteen.
\item The system shall allow students to submit a complaint about a canteen.
\item The system shall allow canteen owners to log in securely and access an owner dashboard.
\item The system shall allow canteen owners to add new food items to their canteen's menu.
\item The system shall allow canteen owners to edit the details of an existing food item, including its price.
\item The system shall allow canteen owners to delete a food item from their canteen's menu.
\item The system shall allow canteen owners to update the availability status of a food item.
\item The system shall allow canteen owners to view all reviews and complaints submitted for their canteen.
\end{itemize}


\subsection{Non- Functional Requirements}
\begin{itemize}
\item \textbf{Usability:} The interface shall be simple and intuitive enough for students and canteen owners with minimal technical background to use without additional training.
\item \textbf{Performance:} The system shall load canteen and food listings within a few seconds under normal usage conditions on the local development environment.
\item \textbf{Security:} Student and owner credentials shall be validated during login, and access to owner-specific functions shall be restricted to authenticated owners.
\item \textbf{Reliability:} The system shall correctly and consistently store and retrieve data for canteens, food items, reviews, and complaints without data loss.
\item \textbf{Maintainability:} The codebase shall be organized so that new features, such as online ordering, can be added in the future with minimal changes to the existing structure.
\item \textbf{Portability:} The system shall run on any environment that supports Apache, PHP, and MySQL, such as XAMPP, without requiring platform-specific modification.
\end{itemize}

\clearpage

\newpage

\section{Results \& Analysis}

The completed CampusBite system was tested on the local XAMPP environment to verify that both the student and canteen owner workflows function as intended. Test student accounts were used to register, log in, browse canteens and their menus, and submit sample reviews and complaints. Test owner accounts were used to log in, add and edit food items, update prices and availability, and view the reviews and complaints submitted by students.

During testing, all core functional requirements listed in Section 5.1 were verified to work as expected: students were able to view canteens and menus, check food prices and availability, and submit reviews and complaints; owners were able to add, edit, and delete food items, update prices and availability, and view the feedback submitted for their canteen. The relational structure of the database, with food items, reviews, and complaints each linked to a specific canteen, ensured that owners only saw information relevant to their own canteen. No major data inconsistencies were observed during repeated create, update, and delete operations on food items during testing.

Overall, the results indicate that CampusBite successfully meets the objectives defined in the scope: it gives students an easy way to find food and give feedback, and gives canteen owners a practical tool to manage their offerings and respond to that feedback.

\clearpage

\newpage

\section{Conclusion \& Future work}

CampusBite was developed to address a simple but persistent problem on university campuses: students have no easy way to find out what food is available, at what price, and where, while canteen owners have no digital tool to manage their offerings or receive structured feedback. By building a web-based system around six core entities -- Student, Owner, Canteen, Food, Review, and Complaint -- and implementing it using HTML, CSS, JavaScript, PHP, and MySQL on a XAMPP environment, this project delivered a working platform that connects students and canteen owners through a shared, centralized source of food and feedback information.

Testing confirmed that the system meets its functional and non-functional requirements, allowing students to browse canteens and menus and submit feedback, and allowing owners to manage food listings and view that feedback through a dedicated dashboard.

The current version of CampusBite is intentionally limited in scope, and several directions for future work follow naturally from the features left out of this version:
\begin{itemize}
    \item Adding online food ordering and digital payment integration so that students can order and pay for food directly through the system.
    \item Introducing food reservation, so that students can reserve an item before arriving at the canteen.
    \item Sending notifications to students about food availability and to owners about new reviews or complaints.
    \item Developing a mobile application version of CampusBite for easier access on students' phones.
    \item Adding a canteen rating system and real-time availability updates.
    \item Building analytics tools to help owners understand popular food items and student preferences.
    \item Introducing an administrative panel for overall system management and moderation of reviews and complaints.
\end{itemize}

These enhancements would extend CampusBite from a food-information and feedback platform into a more complete campus food-service ecosystem.

\clearpage

\newpage

\section{References}
\begin{thebibliography}{9}

\bibitem{b1}
Mozilla Contributors, ``JavaScript,'' \textit{MDN Web Docs}. [Online]. Available: \url{https://developer.mozilla.org/en-US/docs/Web/JavaScript}. [Accessed: 09-Aug-2026].

\bibitem{b2}
Mozilla Contributors, ``HTML: HyperText Markup Language,'' \textit{MDN Web Docs}. [Online]. Available: \url{https://developer.mozilla.org/en-US/docs/Web/HTML}. [Accessed: 09-Aug-2026].

\bibitem{b3}
Mozilla Contributors, ``CSS: Cascading Style Sheets,'' \textit{MDN Web Docs}. [Online]. Available: \url{https://developer.mozilla.org/en-US/docs/Web/CSS}. [Accessed: 09-Aug-2026].

\bibitem{b4}
Oracle Corporation, ``MySQL 8.0 Reference Manual,'' \textit{MySQL Documentation}. [Online]. Available: \url{https://dev.mysql.com/doc/refman/8.0/en/}. [Accessed: 09-Aug-2026].

\bibitem{b5}
Apache Friends, ``XAMPP Documentation,'' \textit{Apache Friends}. [Online]. Available: \url{https://www.apachefriends.org/docs/}. [Accessed: 09-Aug-2026].

\bibitem{b6}
phpMyAdmin Project, ``phpMyAdmin Documentation,'' \textit{phpMyAdmin.net}. [Online]. Available: \url{https://docs.phpmyadmin.net/en/latest/}. [Accessed: 09-Aug-2026].

\end{thebibliography}

\end{document}
